% ============================================================================
% Learning from Big Data (B3MIN1039), Module 2, Assignment 2
% Multi-armed bandits on the ZOZOTOWN Open Bandit data
%
% Submission file name: Assignment_2_STUDENTNUMBER.pdf
% Page limit: 15 pages without the appendix (only the first 15 are graded).
%
% Every number in this report comes from output/run_log.txt, produced by
% `python main.py` in the code repository. Each table and figure caption names
% the table or figure ID printed by the code ("code Table 13").
%
% Parts still to fill in after the final run are marked with \todo{...}.
% ============================================================================

\documentclass[11pt]{article}

\usepackage[a4paper,margin=2.5cm]{geometry}
\usepackage{float}
\usepackage{booktabs}
\usepackage{graphicx}
\usepackage{caption}
\usepackage{amsmath}
\usepackage{amssymb}
\usepackage{xcolor}
\usepackage[round]{natbib}
\usepackage{hyperref}
\hypersetup{colorlinks=true, linkcolor=blue, citecolor=blue, urlcolor=blue}

\captionsetup{font=small, labelfont=bf}
\graphicspath{{../output/}}

% Placeholder for parts that still need the final run or a decision
\newcommand{\todo}[1]{\textcolor{red}{[TODO: #1]}}

% Include a figure from the code output, or a visible box if it does not exist yet
\newcommand{\codefigure}[2]{%
  \IfFileExists{../output/#1}{\includegraphics[width=#2\textwidth]{#1}}%
  {\fbox{\parbox{#2\textwidth}{\centering\vspace{2cm}\textcolor{red}{\detokenize{#1} not found: run main.py}\vspace{2cm}}}}%
}

\newcommand{\studentnumber}{STUDENTNUMBER}

\title{Learning from Big Data: Assignment 2\\[2pt]
       \large Recommending fashion items with Thompson Sampling and UCB}
\author{Dongjie (Leo) Wu Qiu, 721207 \studentnumber}
\date{\today}

\begin{document}
\maketitle

% ============================================================================
\section{Introduction}
% ============================================================================

ZOZOTOWN, a large Japanese fashion website, shows every visitor a row of three
recommended items. Which items should go into that row, and in which slot, to
get the most clicks? With dozens of candidate items and click rates well below
one percent, the site cannot simply test every item at length: every
impression spent on a weak item is a click lost. This is a multi-armed bandit
problem. Each item is an arm, a click is the reward, and the policy has to
balance trying uncertain items (exploration) against showing the items that
look best so far (exploitation).

We use a week of logged data from a randomised experiment that ZOZO ran in
November 2019 \citep{saito2021}. We implement Thompson Sampling (TS) and the
Upper Confidence Bound algorithm (UCB) and evaluate them offline with the
replay method of \citet{li2011}. Beyond the comparison asked for in question
(a), we study batch updates (b), tuning (c), user heterogeneity (d), fairness
(e), stationarity (f) and policies for a changing environment (g). Two
extensions are implemented: a minimum exposure guarantee for items, and
discounted TS and sliding window UCB for non-stationary data.

\todo{Two or three sentences with the headline findings, written last.}

% ============================================================================
\section{Data}
% ============================================================================

The main file has 1{,}356{,}670 impressions of 79 items over seven days
(24 to 30 November 2019), with 4{,}720 clicks, an overall click-through rate
(CTR) of 0.348\%. Each row records one item shown in one slot (1 = left,
2 = centre, 3 = right), whether it was clicked, its propensity score and four
hashed user features. Two smaller files contain the same log restricted to 10
items (172{,}400 rows, CTR 0.280\%) and 20 items (345{,}469 rows, CTR 0.255\%);
the 10 item file has no user features. We use the 80 item file for all main
results and the 10 item file as an easier setting with fewer arms
(code Table~1).

Three facts about the data drive the rest of the analysis.

\paragraph{The logging policy was uniform random.}
Every row has propensity $0.0125 = 1/80$ (code Table~2). The brief mentions
that ZOZO also ran a Bernoulli TS policy, but those impressions are not in this
file, so the policy that generated our data, and our benchmark, is uniform
random serving. Uniform logging is also what makes unbiased replay
evaluation possible (Section~\ref{sec:replay}).

\paragraph{Items differ, positions hardly do.}
Item CTRs range from 0.115\% to 0.784\% (Figure~\ref{fig:itemctr}), so the
best item is clicked almost seven times as often as the worst. That spread is
what a bandit can exploit. The three positions have almost the same CTR:
0.354\%, 0.349\% and 0.341\% (code Table~3).

\begin{figure}[H]
  \centering
  \codefigure{figure_1_item_ctr.png}{0.85}
  \caption{CTR per item in the 80 item file with 95\% Wilson intervals,
    sorted from best to worst; the dashed line is the overall CTR
    (code Figure~1).}
  \label{fig:itemctr}
\end{figure}

\paragraph{One user code covers a third of the data.}
In all four user features the same code covers exactly 479{,}757 rows
(35\%), which suggests these are visitors without a known profile
(code Table~5). Several feature levels are tiny (some have fewer than 100
rows). This matters for questions (d) and (e).

Daily CTR moves between 0.303\% and 0.382\% (code Table~6, Figure~2); we test
in Section~\ref{sec:f} whether this is more than noise.

% ============================================================================
\section{Method}
% ============================================================================

\subsection{Offline evaluation by replay}
\label{sec:replay}

We cannot show items to real visitors, so every policy is evaluated on the
log with the replay method \citep{li2011}. Let the log be a sequence of
events $(x_t, a_t, r_t)$, $t = 1,\dots,T$, in time order, where $x_t$ is the
context (slot and user features), $a_t$ the item that was shown and
$r_t \in \{0,1\}$ the click. A policy $\pi$ walks through the log. At each
event it picks an item $\pi(x_t, h_{t-1})$ based on its history $h_{t-1}$. If
the pick equals the logged item, the event is \emph{matched}: the click is
revealed, counted, and added to the history. Otherwise the event is skipped
and the policy learns nothing. The replay estimate of the policy's CTR is
%
\begin{equation}
  \widehat{\text{CTR}}(\pi)
  = \frac{\sum_{t=1}^{T} r_t \,\mathbf{1}\{\pi(x_t,h_{t-1}) = a_t\}}
         {\sum_{t=1}^{T} \mathbf{1}\{\pi(x_t,h_{t-1}) = a_t\}} .
  \label{eq:replay}
\end{equation}
%
\citet{li2011} show that this estimate is unbiased when the logged item was
chosen uniformly at random and independently of the click, and the events are
independent. The first condition holds by design. With $K$ arms, about one
event in $K$ matches, so both the 10 and the 80 item file leave roughly
17{,}000 matched events and only 45 to 130 clicks per run. One standard error
of a CTR estimate is 0.04 to 0.07 percentage points, and we treat differences
of that size as noise.

The independence condition holds only approximately. In the log, the same
item often stays in the same slot for many consecutive rows. Deterministic
policies such as UCB keep the same pick until they learn something new, so
they match one row of such a streak and then wait for the next streak of
their new favourite: UCB with $c=1$ matches about 5{,}000 events in time order
but about 17{,}000 when the rows are shuffled (code Table~18). We report
matched events next to every CTR and return to this in Section~\ref{sec:c}.

Each stochastic policy is replayed 10 times (5 times in the larger grids),
run $i$ using random seed $1+i$. Tables report the mean CTR over runs, the
spread across runs, and the binomial standard error
$\sqrt{\hat p(1-\hat p)/n}$ of one run with $n$ matched events. The benchmark
at the top is the \emph{Oracle}, which always shows the item with the highest
CTR over the whole week. It uses information no real policy has, so it is an
upper bound for any policy that shows the same item to everyone.

\subsection{Thompson Sampling}

TS keeps a Beta belief about the click probability $\theta_k$ of every item
$k$. Starting from the prior $\text{Beta}(a_0, b_0)$, after $s_k$ clicks and
$f_k$ non-clicks the posterior is
%
\begin{equation}
  \theta_k \mid \text{data} \sim \text{Beta}(a_0 + s_k,\; b_0 + f_k),
  \label{eq:ts}
\end{equation}
%
by Beta-Bernoulli conjugacy. At every event TS draws one value
$\tilde\theta_k$ from each posterior and shows $\arg\max_k \tilde\theta_k$;
after a matched event it adds one to $s_k$ or $f_k$ \citep{russo2018}.
Items with an uncertain posterior sometimes draw a high value and get tried,
so exploration comes from the posterior itself. The prior can be read as
$a_0-1$ imagined clicks and $b_0-1$ imagined non-clicks with mean
$a_0/(a_0+b_0)$. The default $\text{Beta}(1,1)$ is flat over $[0,1]$.

\subsection{Upper Confidence Bound}

UCB1 \citep{auer2002} first shows every item once and then shows the item with
the highest index
%
\begin{equation}
  U_k = \hat\mu_k + c\,\sqrt{\frac{2\ln n}{n_k}},
  \label{eq:ucb}
\end{equation}
%
where $\hat\mu_k$ is the item's CTR so far, $n_k$ the number of times it was
shown and $n$ the total number of matched events. The second term is an
optimism bonus that shrinks as an item is shown more. The textbook version has
$c=1$; we treat $c$ as a tuning parameter. Unlike TS, UCB is deterministic:
the same history always gives the same choice.

\subsection{Measures}

Besides the CTR we report the lift over Random,
$100(\widehat{\text{CTR}}/\widehat{\text{CTR}}_{\text{Random}} - 1)$, the
regret per 10{,}000 matched events,
$10^4(\widehat{\text{CTR}}_{\text{Oracle}} - \widehat{\text{CTR}})$, which
counts the clicks lost against always showing the best item, and a two-sided
two-proportion $z$-test against Random based on the clicks and matched events
of one run.

% ============================================================================
\section{Question (a): comparing the policies}
\label{sec:a}
% ============================================================================

\subsection{Main comparison}

Table~\ref{tab:a} compares the policies on both files. The defaults are TS
with a flat prior and UCB with $c=1$. The tuned versions use the settings
found in question (c): a prior $\text{Beta}(1,300)$, whose mean of 0.33\% is
close to the overall CTR, and $c=0.01$.

\begin{table}[H]
  \centering
  \caption{Replay CTR of the policies, 10 runs each (code Table~13). Lift is
    relative to Random, regret is the number of clicks lost against the
    Oracle per 10{,}000 matched events, $p$ is the two-proportion $z$-test
    against Random. Tuned settings come from question (c) and were chosen on
    the same log.}
  \label{tab:a}
  \footnotesize
  \setlength{\tabcolsep}{4pt}
  \begin{tabular}{lrrrrr@{\hspace{14pt}}rrr}
    \toprule
    & \multicolumn{5}{c}{80 items} & \multicolumn{3}{c}{10 items} \\
    \cmidrule(lr){2-6}\cmidrule(lr){7-9}
    Policy & Matched & CTR (\%) & Lift (\%) & Regret & $p$ & CTR (\%) & Lift (\%) & $p$ \\
    \midrule
    Random (logging policy)        & 17{,}192 & 0.330 &   0.0 & 45.4 &       & 0.261 &   0.0 &       \\
    TS, $\text{Beta}(1,1)$         & 17{,}109 & 0.377 &  14.1 & 40.7 & 0.47  & 0.380 &  45.4 & 0.05  \\
    TS, $\text{Beta}(1,300)$       & 17{,}077 & 0.413 &  25.0 & 37.1 & 0.21  & 0.395 &  51.0 & 0.03  \\
    UCB, $c=1$                     & 10{,}757 & 0.408 &  23.4 & 37.6 & 0.30  & 0.330 &  26.3 & 0.43  \\
    UCB, $c=0.01$                  & 11{,}169 & 0.502 &  52.0 & 28.2 & 0.02  & 0.398 &  52.3 & 0.03  \\
    Per group TS (Section~\ref{sec:d}) & 17{,}073 & 0.443 & 34.2 & 34.1 & 0.09  &       &       &       \\
    Oracle (best item)             & 16{,}325 & 0.784 & 137.3 &  0.0 &       & 0.525 & 100.9 &       \\
    \bottomrule
  \end{tabular}
\end{table}

Random replays at 0.330\% on the 80 item file, close to the 0.348\% of the
full log, and the Oracle replays at exactly the CTR of the best item. Both
checks confirm that the evaluator works (code Table~7).

The two files tell different stories. With 10 items every learning policy
beats Random by 45 to 52\%, and TS shows the best item in 31.5\% of its
matched events against 10\% under Random (code Table~9). With 79 items the
same 17{,}000 matched events are spread over eight times as many arms, about
217 per item and on average less than one click each. Default TS then barely
moves away from uniform: it shows the best item 1.7\% of the time against
1.27\% for Random, and its 14\% lift is not significant ($p=0.47$). The flat
prior is the reason. After $n$ non-clicks its posterior mean is $1/(n+2)$, so
after 217 shows an item without clicks still looks like 0.46\%, above the
true CTR of most items, and TS keeps exploring until the data runs out.
Figure~\ref{fig:learning} shows the result: the TS curve separates from Random
early and then stays flat, far below the Oracle.

\begin{figure}[H]
  \centering
  \codefigure{figure_3_learning_80items.png}{0.75}
  \caption{Cumulative CTR during replay on the 80 item file, mean over runs,
    for TS, UCB ($c=1$) and Random; the dashed line is the Oracle
    (code Figure~3).}
  \label{fig:learning}
\end{figure}

Tuning changes the ranking. The informed prior lifts TS to 0.413\%, and UCB
with $c=0.01$ reaches 0.502\%, a 52\% lift and the only significant result on
the 80 item file ($p=0.02$). With $c=1$ the bonus in~\eqref{eq:ucb} is about
0.2 after a few thousand events, around 70 times the size of the CTRs it is
added to, so UCB mostly rotates through the items; on the 10 item file it shows
every item about 10\% of the time, exactly like Random (code Table~9). Even
the best policy leaves about 60\% of the Oracle's advantage over Random on the
table: there are simply too few clicks to identify the best of 79 items within
one week of uniformly logged data.

\subsection{Where to place the items}

The brief also asks where to place the items. If the best item depended on
the slot, a policy should learn a separate ranking for the left, centre and
right position. The five best items per position overlap heavily (items 33,
34, 38, 57 and 59 appear in several slots, code Table~10). A likelihood ratio
test of the logistic model with item and position effects against a model
with a separate rate for every item and position rejects at the textbook
level ($\text{LR}=332.0$, 156 degrees of freedom, $p<10^{-14}$). That test,
however, assumes independent rows. When we swap the three position labels
between slots at random within every hour, which keeps the streaks intact but
removes any real link between slot and item, the placebo statistics average
323.3 and the permutation $p$-value is 0.38 (code Table~11). We find no
evidence that the best item depends on the slot. Consistently, TS with one
model per position does not beat one model for all positions (0.399\% and
0.395\% against 0.387\%, code Table~12). The practical answer is to show the
best items and not to spend data on learning a ranking per slot.

% ============================================================================
\section{Question (b): batch size}
\label{sec:b}
% ============================================================================

A production system does not update after every impression; it collects
results and updates in batches. With batch size $B$ the policy keeps its
beliefs fixed for $B$ matched events and then processes all of them at once.
We vary $B$ from 1 to 5{,}000 for TS and for UCB with $c=0.01$
(Table~\ref{tab:b}).

\begin{table}[H]
  \centering
  \caption{CTR (\%) for different batch sizes, 5 runs each (code Table~14).
    Random replays at 0.261\% (10 items) and 0.330\% (80 items).}
  \label{tab:b}
  \small
  \begin{tabular}{lrrrrrr}
    \toprule
    Batch size & 1 & 10 & 100 & 500 & 1{,}000 & 5{,}000 \\
    \midrule
    TS, 10 items            & 0.405 & 0.392 & 0.343 & 0.337 & 0.365 & 0.353 \\
    UCB $c=0.01$, 10 items  & 0.425 & 0.403 & 0.421 & 0.348 & 0.368 & 0.274 \\
    TS, 80 items            & 0.387 & 0.353 & 0.370 & 0.316 & 0.356 & 0.366 \\
    UCB $c=0.01$, 80 items  & 0.569 & 0.376 & 0.471 & 0.465 & 0.333 & 0.366 \\
    \bottomrule
  \end{tabular}
\end{table}

On the 10 item file both policies lose clicks as batches grow, but in
different ways. TS loses some of its advantage and then levels off; at
$B=5{,}000$, which leaves only three or four updates in the whole week, it
still gets 0.353\%, well above Random. UCB hardly suffers up to $B=100$ and
then falls to 0.274\%, Random level. The difference follows from the
algorithms. During a batch UCB's index does not change, so it shows the same
item for the whole batch; a weak choice is repeated thousands of times. TS
keeps drawing from its posterior and spreads the batch over several promising
items. This agrees with \citet{chapelle2011}, who found that TS copes well with delayed
feedback. On the 80 item file the picture is weaker: TS barely learns at any
batch size, so there is little to lose, and UCB's results vary so much
between runs (standard deviation across runs up to 0.13 percentage points)
that the drop is not clear beyond $B=1$. A second effect shows in the
matched counts: while UCB holds one choice during a batch it also catches the
streaks in the log, and its matched events rise from 12{,}930 at $B=1$ to
about 17{,}000 at $B\geq 100$ (10 items).

% ============================================================================
\section{Question (c): tuning}
\label{sec:c}
% ============================================================================

We vary the prior of TS and the bonus scale $c$ of UCB (Figure~\ref{fig:tuning}).

\begin{figure}[H]
  \centering
  \codefigure{figure_5_tuning_80items.png}{0.9}
  \caption{CTR for TS with prior $\text{Beta}(1,b)$ (left) and UCB with bonus
    scale $c$ (right) on the 80 item file, 5 runs per setting, with
    $\pm 1.96$ binomial standard errors; the dashed line is Random
    (code Figure~5, code Tables~15 and 16).}
  \label{fig:tuning}
\end{figure}

For TS the prior matters most when there are many arms. On the 80 item file
CTR rises from 0.387\% with the flat prior to 0.420\% with $\text{Beta}(1,300)$
and 0.483\% with $\text{Beta}(1,1000)$, whose prior means of 0.33\% and 0.10\%
are on the scale of the real CTRs. A flat prior claims that a CTR of 50\% is
as likely as one of 0.3\%, and the data spends most of its clicks correcting
that claim. On the 10 item file, with about 1{,}700 matched events per item,
the data overrides any prior and all settings lie within noise of each other
(0.31 to 0.42\%, code Table~15).

UCB is more sensitive. On the 10 item file CTR follows a hill in $c$: 0.300\%
at $c=1$, 0.444\% at $c=0.03$ and 0.341\% at $c=0.001$. A large bonus keeps
UCB exploring; a very small one makes it commit to whichever item looked best
early, which is sometimes an item that was lucky. On the 80 item file the best
values are $c=0.01$ (0.569\%) and $c=0.1$ (0.528\%) (code Table~16). A bonus
on the scale of the CTRs, not on the scale of $[0,1]$ rewards for which UCB1 was
designed, is what makes UCB work here.

These best settings were chosen on the same log on which they are scored, so
they flatter the tuned policies. As a fairer check we choose the setting on
days 1 and 2 only and score it on days 3 to 7 (code Table~17). On the 80 item
file the choice for UCB is $c=0.1$, which scores 0.502\% on the later days
against 0.424\% for the default and 0.378\% for Random, so tuning UCB helps out
of sample. The choice for TS, $\text{Beta}(1,30)$, does not help (0.372\%
against 0.381\%). On the 10 item file the first two days leave only about a dozen
matched clicks, too few to separate the settings, and the defaults are chosen. In
practice $c$ and the prior have to be set before launch, ideally from earlier
data on the site's typical CTR.

Finally, replaying the log in shuffled order removes the streaks. UCB with
$c=1$ then matches 17{,}112 events instead of 10{,}639 on the 80 item file
(code Table~18). The CTRs are lower in shuffled order for five of the six
policy and file combinations, by 0.05 to 0.11 percentage points (for example
0.569\% against 0.501\% for UCB with $c=0.01$). Each gap is about one standard
error, but the consistent direction suggests that replay on this log slightly
favours policies that stay with one item, perhaps because clicks cluster
within streaks. \todo{Random in time order vs shuffled from code Table~18:
if Random shows no gap, write that the gap appears only for adaptive
policies.} The ranking of the policies does not change, but absolute CTRs
should be read with this caveat.

% ============================================================================
\section{Question (d): one model or one per user group}
\label{sec:d}
% ============================================================================

So far every visitor gets the same recommendation. If user groups prefer
different items, one model per group could do better, but each group model
then learns from a fraction of the data. We compare three designs for each of
the four user features: one pooled TS, one TS per level of the feature
starting from $\text{Beta}(1,1)$, and one TS per level that borrows from the
pooled model. In the last design the prior of item $k$ for group $g$ is
%
\begin{equation}
  a_{0,k} = 1 + w_k\,\bar p_k, \qquad b_{0,k} = 1 + w_k\,(1-\bar p_k),
  \qquad w_k = \min(m, N_k),
  \label{eq:borrow}
\end{equation}
%
where $\bar p_k$ and $N_k$ are the CTR and the number of matched events of
item $k$ in the pooled model so far and $m=1{,}000$. A small group stays close
to the pooled estimate; a large group soon outweighs the borrowed $m$ shows
with its own data. This is shrinkage towards the pooled mean. The pooled
counts are built during the replay from past events only, so the policy uses
no information that a live system would not have. A group's own clicks are
counted in both terms, which double counts them slightly; the cap $m$ keeps
this small.

\begin{table}[H]
  \centering
  \caption{CTR (\%) of pooled and per group TS on the 80 item file, 5 runs
    each (code Table~19). Random gives 0.341\% and pooled TS 0.387\% in the
    same runs; one standard error is about 0.05.}
  \label{tab:d}
  \small
  \begin{tabular}{lrrrr}
    \toprule
    & user\_feature\_0 & user\_feature\_1 & user\_feature\_2 & user\_feature\_3 \\
    Levels & 4 & 6 & 10 & 10 \\
    \midrule
    Per group TS, from zero    & 0.342 & 0.374 & 0.357 & 0.405 \\
    Per group TS, borrowing    & 0.494 & 0.388 & 0.362 & 0.395 \\
    \bottomrule
  \end{tabular}
\end{table}

Splitting from zero never beats the pooled model (Table~\ref{tab:d}). Each
group model starts with no knowledge and has too few clicks to learn.
Borrowing helps clearly only for \texttt{user\_feature\_0}, the feature with
the fewest and largest groups: 0.494\% against 0.387\% for the pooled model.
For the features with ten levels, many of them small, it makes no
difference. Within \texttt{user\_feature\_0} all three evaluable groups gain
(code Table~20). Group 3, with 150{,}691 rows, is the clearest case: the
pooled model lifts its CTR by 14\% over Random, the per group model with
borrowing by 100\%, so this group prefers different items from the majority.
Group 0 has 6{,}641 rows and about 80 matched events per run, too few to
learn anything or even to measure how well it is served. Aggregation wastes
the differences between groups; full disaggregation wastes data. Borrowing
gets most of both, but only where groups are large enough.

% ============================================================================
\section{Question (e): fairness}
\label{sec:e}
% ============================================================================

A policy that maximises clicks can be unfair in two ways. It can concentrate
exposure on a few items, so that the other sellers are hardly seen and an item
that was unlucky early never gets a second chance (fairness to items). It can
also serve some user groups much better than others, usually small groups or
groups whose taste differs from the majority (fairness to users).

We measure fairness to items with the Gini coefficient of the exposure shares
$e_k$ (the share of matched events in which item $k$ was shown),
%
\begin{equation}
  G = \frac{2\sum_{i=1}^{K} i\, e_{(i)}}{K\sum_{i} e_{(i)}} - \frac{K+1}{K},
  \label{eq:gini}
\end{equation}
%
with $e_{(1)}\leq\dots\leq e_{(K)}$, so $G=0$ means equal exposure. We also
count the items that receive less than half of the fair share $1/K$, and
report the rank correlation between exposure and item CTR, which shows
whether exposure follows merit. For fairness to users we compute, per level
of \texttt{user\_feature\_0}, the CTR divided by Random's CTR in the same
group (the lift), and report the worst group's lift and the gap between the
best and the worst group. Groups with fewer than 500 matched events per run
are left out because their lift cannot be estimated; on the 80 item file this
excludes group 0.

As fairness extension we add a minimum exposure guarantee. A wrapper around
any policy promises every item at least a fraction $\phi$ of its fair share:
whenever some item has been shown fewer than $\phi\, n/K$ times after $n$
matched events, the wrapper shows the most neglected item; otherwise the base
policy decides. We put the floor on UCB, the policy that concentrates
exposure, with $\phi = 0.25$ and $\phi = 0.5$.

\begin{table}[H]
  \centering
  \caption{Clicks against fairness on the 80 item file, 5 runs each
    (code Table~21). Fair share per item is $1/79 = 1.27\%$.}
  \label{tab:e}
  \footnotesize
  \setlength{\tabcolsep}{4pt}
  \begin{tabular}{lrrrrrr}
    \toprule
    & & \multicolumn{3}{c}{Fairness to items} & \multicolumn{2}{c}{Fairness to users} \\
    \cmidrule(lr){3-5}\cmidrule(lr){6-7}
    Policy & CTR (\%) & Gini & Items $<\tfrac12$ share & Exposure vs CTR & Worst lift & Gap \\
    \midrule
    Random                      & 0.34 & 0.021 &  0 & $-0.06$ & 1.00 & 0.00 \\
    TS                          & 0.39 & 0.090 &  0 & 0.69 & 1.10 & 0.05 \\
    UCB $c=0.01$                & 0.57 & 0.673 & 54 & 0.24 & 0.92 & 1.06 \\
    UCB + floor $\phi=0.25$     & \todo{} & & & & & \\
    UCB + floor $\phi=0.5$      & \todo{} & & & & & \\
    Per group TS, borrowing     & 0.49 & 0.157 &  0 & 0.67 & 1.28 & 0.72 \\
    \bottomrule
  \end{tabular}
\end{table}

The policies sit at different points of the trade-off (Table~\ref{tab:e},
Figure~\ref{fig:fair}). UCB with $c=0.01$ earns the most clicks, but it puts
41\% of all exposure on three items and leaves 54 of the 79 items below half
their fair share. It is also the least fair to users: it lifts group 2 by
97\% and leaves group 3 slightly below Random (lift 0.92, code Table~22),
because its few favourite items suit the majority. TS is close to fair on
both counts: its random draws spread exposure ($G=0.09$), exposure follows
merit (rank correlation 0.69), and every group gains 10 to 15\%. The price is
a lower CTR. Per group TS with borrowing gives the best deal for users, the
highest worst-group lift (1.28) with a moderate Gini of 0.16, because it
serves each large group with its own favourites.

\todo{Floor results from code Table~21 and 22: how much the Gini and the number
of starved items fall for $\phi=0.25$ and $\phi=0.5$, what it costs in CTR,
and how the group lifts change. In the 20 item test the floor $\phi=0.5$ cut
the Gini of UCB from 0.64 to 0.38 and left no item below half its share.}

\begin{figure}[H]
  \centering
  \codefigure{figure_6_fairness_80items.png}{0.7}
  \caption{CTR against the Gini of item exposure on the 80 item file; error
    bars are $\pm 1.96$ binomial standard errors (code Figure~6).}
  \label{fig:fair}
\end{figure}

% ============================================================================
\section{Question (f): is the data stationary?}
\label{sec:f}
% ============================================================================

Two kinds of change matter differently for a bandit. If all items become more
or less popular together, the best item stays the best and a stationary policy
still picks correctly. If the ranking of the items changes, a policy that
averages over the whole past keeps showing yesterday's winner. We test both on
the 80 item file.

All textbook tests assume independent rows, and the streaks in the log violate
that. We therefore calibrate every test with a block permutation placebo: the
168 hours of the week are dealt out at random into placebo days (or halves),
keeping each hour intact, and the test statistic is recomputed 200 times. A
placebo split keeps the dependence between nearby rows but contains no real
trend, so it shows how large the statistic gets by chance. The permutation
$p$-value is the share of placebo statistics at least as extreme as the real
one.

\begin{table}[H]
  \centering
  \caption{Stationarity tests on the 80 item file (code Tables~23 to 25).
    Placebo values come from 200 hour block permutations.}
  \label{tab:f}
  \footnotesize
  \setlength{\tabcolsep}{4pt}
  \begin{tabular}{lrrrrr}
    \toprule
    Test & Statistic & Textbook $p$ & Placebo mean & Placebo 95\% & Perm.\ $p$ \\
    \midrule
    Equal daily CTR, $\chi^2$ (6 df)          & 27.5  & 0.0001          & 12.4  & 24.4   & 0.015 \\
    Item ranking, halves (Spearman $\rho$)    & 0.772 &                 & 0.747 & 0.684$^{a}$ & 0.74 \\
    Item $\times$ half, LR (78 df)            & 136.8 & 0.00004         & 166.4 & 207.1  & 0.89 \\
    Item $\times$ day, LR (468 df)            & 928.7 & $<10^{-6}$      & 972.4 & 1{,}057.8 & 0.82 \\
    \bottomrule
    \multicolumn{6}{l}{\footnotesize $^{a}$ 5th percentile, since a change in the ranking lowers $\rho$.}
  \end{tabular}
\end{table}

The textbook tests all reject stationarity, but the placebo values show why
they cannot be trusted here: even random groupings of hours give a $\chi^2$ of
12.4 on average instead of the 6 expected under independence
(Table~\ref{tab:f}). After calibration the results split. The overall daily
CTR does change: the real $\chi^2$ of 27.5 beats 98.5\% of the placebo values
(permutation $p=0.015$), with Wednesday at 0.303\% and Tuesday at 0.382\%. The
item ranking does not. The rank correlation of item CTRs between days 1 to 3
and days 4 to 7 is 0.772, a typical value for a placebo split of hours (mean
0.747, $p=0.74$; Figure~\ref{fig:halves}). The likelihood ratio tests of an
item by period interaction in a logistic model of the click give statistics
below the placebo mean ($p=0.89$ by half week and $p=0.82$ by day). In the
day-level test 23\% of the cells have fewer than five expected clicks, which
is one more reason to rely on the permutation version.

\begin{figure}[H]
  \centering
  \codefigure{figure_7_half_split_80items.png}{0.48}
  \caption{CTR of every item in days 1 to 3 against days 4 to 7, 80 item
    file; points on the diagonal behave the same in both halves
    (code Figure~7).}
  \label{fig:halves}
\end{figure}

The data is therefore non-stationary in level but stationary in ranking. The
level shift does not change which item is best, so it does not hurt a bandit
that only needs the ranking.

% ============================================================================
\section{Question (g) and extension: policies for a changing environment}
\label{sec:g}
% ============================================================================

If the ranking did change, the policies would have to stop weighting old and
new evidence equally. We implement two such changes and outline a third.

Discounted TS \citep{raj2017} multiplies all counts by $\gamma<1$ after every
matched event before adding the new result,
%
\begin{equation}
  s_k \leftarrow \gamma s_k + r\,\mathbf{1}\{k=a\}, \qquad
  f_k \leftarrow \gamma f_k + (1-r)\,\mathbf{1}\{k=a\},
  \label{eq:dts}
\end{equation}
%
and samples from $\text{Beta}(a_0+s_k, b_0+f_k)$. An event that is $n$ updates
old keeps weight $\gamma^n$, a memory of about $1/(1-\gamma)$ updates; we use
$\gamma \in \{0.9999, 0.999, 0.995\}$. Sliding window UCB
\citep{garivier2011} computes $\hat\mu_k$ and $n_k$ in~\eqref{eq:ucb} from the
last $W$ matched events only, with $W \in \{5{,}000, 2{,}000, 1{,}000\}$ and
$c=0.01$. The third option is change point detection, which monitors each
item's recent CTR and restarts the policy when it departs from the long-run
value. With 45 to 130 clicks per run spread over 10 to 79 items, a detector
would have far too little data per item to tell a real change from noise, so
we do not implement it.

Because the ranking in our data is stable (Section~\ref{sec:f}), we also build
a log in which it is not: from the middle of the week the item labels are
reversed by rank, so the rows of the worst item now carry the label of the
best item, the second worst that of the second best, and so on. The logging
stays uniform, so replay remains valid. The cost of the flip is a policy's CTR
in the second half of the real log minus its CTR in the second half of the
flipped log.

\begin{table}[H]
  \centering
  \caption{Forgetting policies on the real log and on the log with the ranking
    flipped halfway, 5 runs each (code Tables~26 and 27). Cost of the flip in
    percentage points; lower means the policy adapted better.}
  \label{tab:g}
  \small
  \begin{tabular}{lrrrr}
    \toprule
    & \multicolumn{2}{c}{10 items} & \multicolumn{2}{c}{80 items} \\
    \cmidrule(lr){2-3}\cmidrule(lr){4-5}
    Policy & CTR real log (\%) & Cost of flip & CTR real log (\%) & Cost of flip \\
    \midrule
    TS (never forgets)                 & 0.405 &  0.191 & 0.387 &  0.048 \\
    Discounted TS $\gamma=0.9999$      & 0.298 & $-0.001$ & 0.374 & $-0.066$ \\
    Discounted TS $\gamma=0.999$       & 0.334 &  0.035 & 0.391 & $-0.001$ \\
    Discounted TS $\gamma=0.995$       & 0.330 &  0.012 & 0.351 &  0.022 \\
    UCB $c=0.01$ (never forgets)       & 0.425 &  0.153 & 0.569 &  0.103 \\
    Sliding window UCB $W=5{,}000$     & 0.457 &  0.071 & 0.567 &  0.106 \\
    Sliding window UCB $W=2{,}000$     & 0.409 &  0.054 & 0.513 &  0.053 \\
    Sliding window UCB $W=1{,}000$     & 0.361 & $-0.062$ & 0.482 &  0.056 \\
    \bottomrule
  \end{tabular}
\end{table}

On the real log, forgetting costs clicks or at best changes nothing
(Table~\ref{tab:g}). On the 10 item file discounted TS loses 0.07 to 0.11
percentage points against plain TS, about two standard errors, and the
shortest window costs UCB 0.06. Throwing away old clicks does not pay when
the ranking is stable and clicks are scarce. On the flipped log the forgetting
policies adapt better. The flip costs plain TS 0.19 percentage points on the
10 item file and every discounted version at most 0.04. On the same file
sliding window UCB pays between nothing and half of the 0.15 that plain UCB
pays. Figure~\ref{fig:shift} shows the mechanism: plain TS drops to about
0.1\% after the flip and needs around 15{,}000 rows to recover, while the
forgetting policies dip less and recover sooner. On the 80 item file the
effects are smaller, since TS had learned little before the flip and has
little to unlearn. Discounted TS with $\gamma=0.9999$ is the only version
with a noticeable gain there, and the longest window does not help UCB at
all (cost 0.106 against 0.103).

\begin{figure}[H]
  \centering
  \codefigure{figure_8_shift_10items.png}{0.75}
  \caption{Moving CTR over the last 1{,}000 matched events on the 10 item log
    with the ranking flipped halfway (dashed line), mean over runs
    (code Figure~8). The peak of sliding window UCB rests on about ten clicks
    per window and overstates its level.}
  \label{fig:shift}
\end{figure}

The choice of model therefore depends on the type of change. On the ZOZO
data, with a level shift and a stable ranking, plain TS or UCB is the better
choice. Forgetting is the right tool only when evidence suggests that the
ranking moves, and its memory length should then be set as long as the speed
of change allows.

% ============================================================================
\section{Limitations and conclusion}
% ============================================================================

\todo{Limitations, one paragraph: (1) few clicks after replay, so most
differences are within one or two standard errors; (2) replay assumes
independent events, the streaks break this and appear to favour sticky
policies slightly (code Table~18); (3) tuned settings chosen in hindsight,
partly checked out of sample; (4) one week of data, Black Friday inside it;
(5) single-slot evaluation, the three slots are not evaluated as a set,
because the log does not record which three items were shown together.}

\todo{Conclusion, one paragraph: answer the business question (which items,
which slot, which policy) in plain terms, with the main numbers from
Table~\ref{tab:a}.}

% ============================================================================
\begin{thebibliography}{9}
\setlength{\itemsep}{2pt}
\small

\bibitem[Auer et al.(2002)]{auer2002}
Auer, P., Cesa-Bianchi, N., and Fischer, P. (2002). Finite-time analysis of the
multiarmed bandit problem. \emph{Machine Learning}, 47(2), 235-256.

\bibitem[Chapelle and Li(2011)]{chapelle2011}
Chapelle, O., and Li, L. (2011). An empirical evaluation of Thompson Sampling.
\emph{Advances in Neural Information Processing Systems}, 24.

\bibitem[Garivier and Moulines(2011)]{garivier2011}
Garivier, A., and Moulines, E. (2011). On upper-confidence bound policies for
switching bandit problems. \emph{Algorithmic Learning Theory}, 174-188.

\bibitem[Li et al.(2011)]{li2011}
Li, L., Chu, W., Langford, J., and Wang, X. (2011). Unbiased offline evaluation
of contextual-bandit-based news article recommendation algorithms.
\emph{Proceedings of WSDM}, 297-306.

\bibitem[Raj and Kalyani(2017)]{raj2017}
Raj, V., and Kalyani, S. (2017). Taming non-stationary bandits: a Bayesian
approach. \emph{arXiv preprint} 1707.09727.

\bibitem[Russo et al.(2018)]{russo2018}
Russo, D., Van Roy, B., Kazerouni, A., Osband, I., and Wen, Z. (2018). A
tutorial on Thompson Sampling. \emph{Foundations and Trends in Machine
Learning}, 11(1), 1-96.

\bibitem[Saito et al.(2021)]{saito2021}
Saito, Y., Aihara, S., Matsutani, M., and Narita, Y. (2021). Open Bandit
Dataset and Pipeline: towards realistic and reproducible off-policy
evaluation. \emph{NeurIPS Datasets and Benchmarks Track}.

\end{thebibliography}

% ============================================================================
\clearpage
\appendix
\section{Appendix}
% ============================================================================
% Unlimited pages. Supplementary tables and figures, reproduction notes.
% ============================================================================

\subsection{Reproduction}

All code is in Python and runs with \texttt{python main.py}, which prints every
table and figure with its ID and saves them to \texttt{output/} (tables as csv
and \LaTeX, the full printout in \texttt{output/run\_log.txt}). The random seed
is fixed (\texttt{random.seed(1)}, and run $i$ of every replay uses seed
$1+i$), so a rerun gives identical numbers. The full run takes about three
hours. \todo{Repository link, or note that the code is uploaded with the
report.}

\begin{table}[H]
  \centering
  \caption{Where each question's results are printed by the code.}
  \label{tab:map}
  \small
  \begin{tabular}{lll}
    \toprule
    Question & Code tables & Code figures \\
    \midrule
    Data                      & 1 to 6   & 1, 2 \\
    (a) comparison, placement & 7 to 13  & 3 \\
    (b) batch size            & 14       & 4 \\
    (c) tuning                & 15 to 18 & 5 \\
    (d) user groups           & 19, 20   & \\
    (e) fairness              & 21, 22   & 6 \\
    (f) stationarity          & 23 to 25 & 7 \\
    (g) forgetting policies   & 26, 27   & 8 \\
    \bottomrule
  \end{tabular}
\end{table}

\subsection{Supplementary figures}

\begin{figure}[H]
  \centering
  \codefigure{figure_2_daily_ctr.png}{0.6}
  \caption{Daily CTR on the 80 item file with 95\% Wilson intervals
    (code Figure~2).}
\end{figure}

\begin{figure}[H]
  \centering
  \codefigure{figure_3_learning_10items.png}{0.7}
  \caption{Cumulative CTR during replay on the 10 item file (code Figure~3).}
\end{figure}

\begin{figure}[H]
  \centering
  \codefigure{figure_4_batch_10items.png}{0.65}
  \caption{CTR against batch size on the 10 item file (code Figure~4).}
\end{figure}

\begin{figure}[H]
  \centering
  \codefigure{figure_5_tuning_10items.png}{0.85}
  \caption{Tuning results on the 10 item file (code Figure~5).}
\end{figure}

\subsection{Supplementary tables}

\todo{Optional: per group lifts for all four user features (code Table~20),
the placement tables for the 10 item file (code Tables~10 to 12), and the
first and second half CTRs of the forgetting policies (code Tables~26 and
27). Only add what the main text refers to.}

\end{document}
